\section{De responder a actuar: Agent OS y exigencias de alineación}

Esta sección extiende el problema de la conducción autónoma a los Agents. Li Xiang (李想) considera que en cinco años no necesariamente habrá un Agent general, pero sí habrá un Agent OS. Este juicio significa lo siguiente: cada oficio profesional necesita un Agent distinto, pero todos esos Agents deben operar sobre una base común similar a un sistema operativo, de la cual obtienen herramientas, permisos, contexto, máquinas virtuales, memoria y límites de seguridad. Un Agent OS no es una ventana de chat, sino la infraestructura de los Agents profesionales.

\begin{figure}[H]
\centering
\includegraphics[width=0.96\textwidth]{action-agent-alignment.png}
\caption{De responder a actuar: cuando la IA interviene sobre activos, herramientas y vehículos, las exigencias de alineación aumentan. Diagrama conceptual propio, elaborado a partir de la conversación de 00:53:58--00:55:00 y 00:55:00--00:55:25.}
\end{figure}

\begin{knowledgebox}{Lectura de la figura: la acción amplifica el riesgo}
Si una búsqueda de información falla, se puede buscar de nuevo; si una investigación tiene errores, se puede revisar. Pero si un Agent opera activos, herramientas, computadoras, vehículos o bienes, el error modifica el mundo de forma directa. Por eso las exigencias de alineación, los límites de permisos y la atribución de responsabilidades de un Agent son mayores que las de un modelo conversacional.
\end{knowledgebox}

\subsection{Por qué un Agent general rinde menos que un Agent profesional}

Li Xiang sostiene que las herramientas de producción reales deben sustituir y liberar a las personas de su trabajo de alta frecuencia, como conducir, programar, operar, vender o fabricar. Cada oficio tiene sus propios datos, herramientas, permisos, costos de error y criterios de evaluación; por ello, a corto plazo es más probable que surjan Agents profesionales que un Agent capaz de todo. El valor del Agent OS consiste en ofrecer a esos Agents profesionales un entorno de ejecución común.

\begin{importantbox}{Definición de Agent OS}
Un Agent OS es el sistema base sobre el que operan los Agents profesionales. Incluye, como mínimo, gestión de contexto, invocación de herramientas, control de permisos, máquinas virtuales o entornos de ejecución, memoria, registros, evaluación y políticas de seguridad. Su valor no está en «poder hacer de todo», sino en permitir que distintos Agents profesionales hagan su trabajo de forma confiable.
\end{importantbox}

\subsection{La alineación pasa de ser un problema de lenguaje a un problema de acción}

Cuando la IA solo responde preguntas, la alineación se refiere sobre todo a la seguridad del contenido, la exactitud factual y las preferencias de valores; cuando la IA empieza a invocar herramientas y activos, la alineación se convierte en riesgo operativo. Por ejemplo, el error de un Agent de conducción afecta la seguridad vial; el de un Agent financiero afecta los fondos; el de un Agent de código afecta la estabilidad del sistema. Así, la alineación se amplía de «si el modelo dice lo correcto» a «si lo que el modelo hace es controlable».

\begin{knowledgebox}{Estratificación aproximada de las exigencias de alineación}
\[
\text{risk} \approx \text{action scope} \times \text{asset value} \times \text{irreversibility}.
\]
Aquí, \(\text{action scope}\) indica el alcance de lo que el Agent puede operar, \(\text{asset value}\) indica el valor de los activos invocados e \(\text{irreversibility}\) indica qué tan difícil es deshacer un error. Mientras más cerca se esté de vehículos, dinero, cuentas, sistemas de producción y dispositivos físicos, mayores serán las exigencias de alineación.
\end{knowledgebox}

\subsection{Resumen de la sección}

El Agent OS es el puente que, en este episodio, lleva de la conducción autónoma a los productos de IA de uso general. Muestra que el siguiente paso de la IA no es conversar mejor, sino actuar mejor en entornos con permisos, herramientas y responsabilidades.
